% ch08 译文：原文 §5b–§5c（main.tex 第 2170–2670 行）
\subsection{特征群 $\widehat{G}$ 在 $\Mod_G(X)$ 上的作用}

在 $X$ 上 $G$-等变层构成的 Abel 范畴 $\Mod_G(X)$ 中，有一个特定的 Abel 群 $\widehat{G}$ 的典范作用；$\widehat{G}$ 称为 $G$ 的\emph{特征群}（character group），下面给出定义。

\begin{defn}
% [原文 p.33]
群概形 $(G,m,e)$ 的\emph{特征}（\emph{character}）指从 $G$ 到乘法群概形
$
\Gm := \Spec\left(\field[t, t^{-1}]\right)
$
的群概形同态
$
\cha : G \longrightarrow \Gm
$。
等价地，特征对应于平凡 $\Gm$-丛 $\mathcal{O}_G^\times$ 在 $G$ 上的一个整体截面
$
\cha
\in
\Gamma\left(G, \mathcal{O}_G^\times\right)
$，
满足
$$
m^* \cha
=
\left(\pr_{G,1}^{G \times G}\right)^* \cha
\cdot
\left(\pr_{G,2}^{G \times G}\right)^* \cha
\in
\Gamma\left(G \times G, \mathcal{O}_{G \times G}^\times\right),
$$
其中
$\pr_{G,1}^{G \times G}$、
$\pr_{G,2}^{G \times G}
:
G \times G \rightarrow G
$
分别为到第一、第二因子的投影，
$
\left(\pr_{G,i}^{G \times G}\right)^* \cha
\in
\Gamma\left(
G \times G,
\left(\pr_{G,i}^{G \times G}\right)^*\mathcal{O}_G^\times
\cong
\mathcal{O}_{G \times G}^\times
\right)
$
是 $\cha$ 沿 $\pr_{G,i}^{G\times G}$ 的拉回（$i\in\{1,2\}$），而 $\cdot$ 表示环
$\mathcal{O}_{G\times G}^\times(G\times G)$ 中元素的乘法。

$G$ 的所有特征构成一个 Abel 群，其群结构由 $\Gm$ 的乘法群概形结构诱导，记作
$$
\widehat{G}
:=
\operatorname{Hom}_{\mathsf{GrpSch}(\field)}\left(G, \Gm\right),
$$
并称为 $G$ 的\emph{特征群}（\emph{character group}）。
\end{defn}

\begin{yizhu}[「特征」的两种等价说法]
特征既可以看成群概形同态 $\cha:G\rightarrow\Gm$，也可以看成 $\Gamma(G,\mathcal{O}_G^\times)$ 中满足 $m^*\cha=(\pr_{G,1})^*\cha\cdot(\pr_{G,2})^*\cha$ 的可逆整体截面：前者由 $\Gm$ 的泛性质与 $\Gm=\Spec(\field[t,t^{-1}])$ 的坐标环决定，后者正是把该同态写成「$t\mapsto\cha$」后的样子。乘法条件保证 $\cha(gh)=\cha(g)\cha(h)$，所以全体特征在逐点相乘下构成 Abel 群 $\widehat{G}$。后文把 $\cha$ 同时当作「群同态」与「可逆截面」使用，正是靠这一等同；$\widehat{G}$ 交换这一事实，则是 §5c 中权空间分解能按特征标逐项拆开的根本原因。

另需说明：原文此处把该丛误写作「the trivial $\Gm$-bundle $\mathcal{O}_X^\times$ on $X$」，但紧接着的公式写的是 $\Gamma(G,\mathcal{O}_G^\times)$ 与 $\mathcal{O}_{G\times G}^\times$，可见所论丛应是 $G$ 上的 $\mathcal{O}_G^\times$；译文按此处理。
\end{yizhu}

\begin{defn}[按特征作扭]
对每个特征
$\cha \in \widehat{G}\subset\Gamma\left(G, \mathcal{O}_G^\times\right)$，
\emph{按 $\cha$ 扭过的平凡线丛}定义为
$$
\mathcal{O}_X(\cha)
:=
(\mathcal{O}_X, \id_{\mathcal{O}_{G\times X}}(\cha))
\in \vb_G(X),
$$
其中底下的 $\mathcal{O}_X$-模是 $\mathcal{O}_X$，$G$-线性化
$$
\id_{\mathcal{O}_{G\times X}}(\cha)
:=
\left(\pr_G^{G \times X}\right)^* \cha \cdot \id_{\mathcal{O}_{G\times X}}
:
\mathcal{O}_{G\times X}
\xrightarrow{\cong}
\mathcal{O}_{G\times X}
$$
（在典范同构
$
\sigma^* \mathcal{O}_X
\cong
\mathcal{O}_{G\times X}
$
与
$
\left(\pr_X^{G \times X}\right)^* \mathcal{O}_X
\cong
\mathcal{O}_{G\times X}
$
之下）由 $\cha$ 沿自然投影
$\pr_G^{G\times X}:G\times X\rightarrow G$ 的拉回
$$
\left(\pr_G^{G \times X}\right)^* \cha
\in
\Gamma\left(G \times X, \left(\pr_G^{G\times X}\right)^*\mathcal{O}_G^\times \cong \mathcal{O}_{G \times X}^\times\right)
$$
与恒等线性化的乘法定义。更确切地说，对每个开集
$U\subseteq G\times X$
与
$s\in\mathcal{O}_{G\times X}(U)$，
定义
$$
\left(\left(\id_{\mathcal{O}_{G\times X}}(\cha)\right)(U)\right)(s)
:=
\left(\left.\left(\pr_G^{G \times X}\right)^* \cha\right|_U\right)
\cdot
s
\in
\mathcal{O}_{G\times X}(U),
$$
其中 $\cdot$ 表示环
$\mathcal{O}_{G\times X}(U)$ 中元素的乘法。

设 $(\mathcal{F}, \theta)$ 为 $\Mod_G(X)$ 中的 $G$-等变 $\mathcal{O}_X$-模，$\cha \in \widehat{G}$。
$(\mathcal{F},\theta)$ \emph{按 $\cha$ 的扭}定义为
$$
(\mathcal{F},\theta)(\cha)
:=
(\mathcal{F}, \theta(\cha))
\in \Mod_G(X),
$$
其中底下的 $\mathcal{O}_X$-模是 $\mathcal{F}$，$G$-线性化
$$
\theta(\cha)
:=
\left(\pr_G^{G \times X}\right)^* \cha \cdot \theta
:
\sigma^* \mathcal{F}
\xrightarrow{\cong}
\left(\pr_X^{G \times X}\right)^* \mathcal{F}
$$
由 $\cha$ 沿自然投影 $\pr_G^{G\times X}$ 的拉回
$$
\left(\pr_G^{G \times X}\right)^* \cha
\in
\Gamma\left(G \times X, \left(\pr_G^{G\times X}\right)^*\mathcal{O}_G^\times \cong \mathcal{O}_{G \times X}^\times\right)
$$
对原 $G$-线性化 $\theta$ 的数乘定义。更确切地说，对每个开集
$U\subseteq G\times X$
与
$s\in(\sigma^*\mathcal{F})(U)$，
定义
$$
(\theta_{\cha}(U))(s)
:=
\left(\left.\left(\pr_G^{G \times X}\right)^* \cha\right|_U\right)
\cdot
(\theta(U)(s))
\in
\left(\left(\pr_X^{G \times X}\right)^*\mathcal{F}\right)(U),
$$
其中 $\cdot$ 表示 $\mathcal{O}_{G\times X}(U)$ 中元素对 $\mathcal{O}_{G\times X}(U)$-模
$\left(\left(\pr_X^{G \times X}\right)^*\mathcal{F}\right)(U)$ 中元素的数乘。

按特征作扭定义了 $\Mod_G(X)$ 上的一个严格 $\widehat{G}$-作用（strict $\widehat{G}$-action）
$$
\phi:\widehat{G}\rightarrow\Aut_{\str}(\Mod_G(X)),
\quad
\cha\mapsto \phi_{\cha},
$$
其中
$$
\phi_{\cha}:\Mod_G(X)\rightarrow\Mod_G(X),
\quad
(\mathcal{F},\theta)
\mapsto
(\mathcal{F},\theta(\cha)).
$$
% [原文 p.34]
这限制为 $\QCoh_G(X)$、$\Coh_G(X)$ 与 $\VB_G(X)$ 上的严格 $\widehat{G}$-作用。
\end{defn}

\begin{rmk}
存在自然同构
$$
(\mathcal{F},\theta)(\cha)
\cong
(\mathcal{F},\theta)
\underset{\mathcal{O}_X}{\otimes}
\mathcal{O}_X(\cha)
\in\Mod_G(X).
$$
\end{rmk}

\subsection{权空间分解}

由于 $\widehat{G}$ 交换，它在 $\Mod_G(X)$ 中同态空间上的作用会把同态空间分解为一维不可约 $\widehat{G}$-表示的直和；我们称之为\emph{权空间分解}（weight space decomposition）。

\begin{conceptbox}{权空间分解}
权空间分解是「交换群作用在向量空间上必可对角化」这一线性代数事实在层论中的推广：当群是对角化群概形 $G=\Spec(\field[\Lambda])$ 时，$\Mod_G(X)$ 中每个对象的整体截面都带 $\field[\widehat{G}]$ 的余模结构（即 $\widehat{G}$-分次结构），于是按特征标 $\cha\in\widehat{G}$ 拆成直和 $\mathcal{F}(X)=\bigoplus_{\cha}\mathcal{F}(X)_{\cha}$，每个 $\mathcal{F}(X)_{\cha}$ 就是「权为 $\cha$ 的部分」。更关键的是，Hom 空间本身也按权分解，且权 $\cha$ 的那一块恰是 $\Hom_{\Mod_G(X)}((\mathcal{F},\theta)(\cha),(\mathcal{F}',\theta'))$ —— 即「把源扭过 $\cha$ 之后」的等变同态。本文 B 侧正是靠这一分解，把 $\mathbb{P}^1$ 上 $\mathbb{G}_m$-等变层的同态空间按权拆开，从而与 A 侧绕 $S^1$ 的单值性、停止位选择逐项对应；§5c 的推论把这一分解提升为轨道范畴的全忠实嵌入，是证明主定理的重要中间步骤，也是 §6 扭 Künneth 公式的技术核心。
\par\smallskip\noindent{\footnotesize 参考：nLab，\emph{equivariant sheaf}；本文 §5b–§5c、§6。}
\end{conceptbox}

\begin{defn}
设 $\Lambda$ 为 Abel 群，$\field$ 为交换环。
\emph{群代数}（\emph{group algebra}）$\field[\Lambda]$ 是一个以 $\{e_\lambda\}_{\lambda \in \Lambda}$ 为基的 Hopf 代数，其中余乘法 $\Delta : \field[\Lambda] \rightarrow \field[\Lambda] \underset{\field}{\otimes} \field[\Lambda]$ 与余单位 $\varepsilon : \field[\Lambda] \rightarrow \field$ 定义为
$$
\Delta(e_\lambda) = e_\lambda \otimes e_\lambda \quad \text{与} \quad \varepsilon(e_\lambda) = 1
\quad\textnormal{对一切 $\lambda \in \Lambda$}.
$$

$\field[\Lambda]$ 上的\emph{右余模}（\emph{right comodule}）指一个 $\field$-模 $V$，连同 $\field$-线性映射 $\rho : V \rightarrow V \underset{\field}{\otimes} \field[\Lambda]$，满足余结合公理与余单位公理：
$$
\rho \comp (\rho \otimes \id_{\field[\Lambda]}) = \rho \comp (\id_V \otimes \Delta)
\quad\text{与}\quad
\rho \comp (\id_V \otimes \varepsilon) = \id_V.
$$
\end{defn}

\begin{prop}[{$\field[\Lambda]$-余模的权空间分解}]
\label{prop:WeightSpaceDecompositionForComodule}
$\field$-模 $V$ 带有 $\field[\Lambda]$ 上的右余模结构，当且仅当 $V$ 是 $\Lambda$-分次的，即 $V$ 分解为权空间的直和
$$
V = \bigoplus_{\lambda \in \Lambda} V_\lambda,
$$
其中 $V_\lambda = \{v \in V \mid \rho(v) = v \otimes e_\lambda\}$。
\end{prop}

\begin{proof}
对任意 $v \in V$，它在 $\rho$ 下的像可唯一写成有限和
$$
\rho(v) = \sum_{\lambda \in \Lambda} P_\lambda(v) \otimes e_\lambda,
$$
其中 $P_\lambda : V \rightarrow V$ 为 $\field$-线性映射。把余单位公理 $\rho \comp (\id_V \otimes \varepsilon) = \id_V$ 作用于 $v$，得到
$$
v = (\id_V \otimes \varepsilon)(\rho(v)) = (\id_V \otimes \varepsilon)\left(\sum_{\lambda \in \Lambda} P_\lambda(v) \otimes e_\lambda\right) = \sum_{\lambda \in \Lambda} P_\lambda(v) \cdot 1 = \sum_{\lambda \in \Lambda} P_\lambda(v).
$$
这说明各 $P_\lambda$ 的像之和覆盖 $V$。其次，把余结合公理
$
\rho \comp (\rho \otimes \id_{\field[\Lambda]}) = \rho \comp (\id_V \otimes \Delta)
$
作用于 $v$，得到
$$
\sum_{\lambda, \mu \in \Lambda} P_\mu(P_\lambda(v)) \otimes e_\mu \otimes e_\lambda = \sum_{\lambda \in \Lambda} P_\lambda(v) \otimes e_\lambda \otimes e_\lambda.
$$
比较两边基元素 $e_\mu \otimes e_\lambda$ 的系数，得 $P_\mu(P_\lambda(v)) = \delta_{\lambda, \mu} P_\lambda(v)$。按从左到右的复合约定，这意味着 $P_\lambda \comp P_\mu = \delta_{\lambda, \mu} P_\lambda$。于是 $\{P_\lambda\}_{\lambda \in \Lambda}$ 是一族相互正交且求和为恒等的投影算子。因此 $V$ 完全分解为像 $\mathrm{Im}(P_\lambda)$ 的直和。容易验证 $\mathrm{Im}(P_\lambda) = V_\lambda$，证毕。
\end{proof}

% [原文 p.35]
\begin{defn}[对角化群概形]
$\field$ 上的仿射群概形 $G$ 称为\emph{对角化}（\emph{diagonalizable}）的，若存在 Abel 群 $\Lambda$ 使得
$G\cong\Spec(\field[\Lambda])$
作为群概形成立。此时，群 $\Lambda$ 与 $G$ 的特征群 $\widehat{G}$ 自然等同，其坐标环 $\mathcal{O}_G(G)$ 作为 Hopf 代数同构于群代数 $\field[\widehat{G}]$。
\end{defn}

\begin{yizhu}[原文笔误：$\Lambda$ 与 $\cha$]
原文此处写作「the group $\cha$ is canonically identified with the character group $\widehat{G}$」，其中的 $\cha$ 显然是笔误，应为上面出现的 Abel 群 $\Lambda$：$\Lambda$ 才是 $\field[\Lambda]$ 的指标群。典范等同来自泛性质：群概形同态 $\Spec(\field[\Lambda])\rightarrow\Gm$ 对应 Hopf 代数同态 $\field[t,t^{-1}]\rightarrow\field[\Lambda]$，而后者由 $t\mapsto e_{\cha}$（$\cha\in\Lambda$）唯一决定；于是 $\widehat{G}$ 自然地重新认出 $\Lambda$ 本身。译文按 $\Lambda$ 处理。
\end{yizhu}

\begin{lemma}[整体截面的权空间分解]
\label{lem:GlobalSectionDecomposition}
设 $G$ 为对角化群概形，作用在概形 $X$ 上，$(\mathcal{F}, \theta) \in \Mod_G(X)$ 为 $G$-等变 $\mathcal{O}_X$-模。则整体截面空间 $V := \mathcal{F}(X)$（$=\Gamma(X, \mathcal{F})$）自然带有 $\mathcal{O}_G(G)$ 上的右余模结构，从而有典范的权空间分解
$$
\mathcal{F}(X) = \bigoplus_{\cha \in \widehat{G}}
\mathcal{F}(X)_{\cha}.
$$
其中
$
\mathcal{F}(X)_{\cha}
:=
\left\{
s \in \mathcal{F}(X)
\left|
(\theta(G\times X))(\sigma^* s)
=
\left(\pr_G^{G \times X}\right)^*\cha
\cdot
\left(\pr_X^{G \times X}\right)^* s
\right.\right\}.
$
\end{lemma}

\begin{proof}
$G$-线性化 $\theta : \sigma^* \mathcal{F} \xrightarrow{\cong} \left(\pr_X^{G \times X}\right)^* \mathcal{F}$ 在整体截面上诱导自然映射。对任意截面 $s \in V = \mathcal{F}(X)$，沿作用映射 $\sigma:G\times X \rightarrow X$ 拉回得到截面 $\sigma^* s \in (\sigma^*\mathcal{F})(G\times X)$。把同构 $\theta$ 作用于该拉回，得
$$
(\theta(G\times X))(\sigma^* s)
\in
\left(\left(\pr_X^{G \times X}\right)^* \mathcal{F}\right)(G\times X)
\cong
\mathcal{F}(X) \underset{\field}{\otimes} \mathcal{O}_G(G)
=
V \underset{\field}{\otimes} \mathcal{O}_G(G).
$$
这定义了 $\field$-线性映射
$$
\rho
:
V \rightarrow V \underset{\field}{\otimes} \mathcal{O}_G(G),
\quad
s
\mapsto
(\theta(G\times X))(\sigma^* s).
$$

$\theta$ 的余循环条件与单位条件分别精确对应于 $\rho$ 的余结合公理与余单位公理。因此 $\rho$ 给出 $V$ 在 $\mathcal{O}_G(G) \cong \field[\widehat{G}]$ 上的右余模结构。由命题 \ref{prop:WeightSpaceDecompositionForComodule}，该余模结构把 $V$ 完全分解为以特征群 $\widehat{G}$ 为下标的权空间直和。

在典范等同
$
\left(\left(\pr_X^{G \times X}\right)^* \mathcal{F}\right)(G\times X)
\cong
\mathcal{F}(X) \underset{\field}{\otimes} \mathcal{O}_G(G)
$
之下，纯张量 $s \otimes e_{\cha} \in V \underset{\field}{\otimes} \mathcal{O}_G(G)$ 对应于
$
\left(\pr_G^{G \times X}\right)^*\cha
\cdot
\left(\pr_X^{G \times X}\right)^* s
\in
\left(\left(\pr_X^{G \times X}\right)^* \mathcal{F}\right)(G\times X)
$，
其中 $\cdot$ 表示环 $\mathcal{O}_{G\times X}(G\times U)$ 中元素对 $\mathcal{O}_{G\times X}(G\times U)$-模
$\left(\left(\pr_X^{G \times X}\right)^*\mathcal{F}\right)(G\times U)$ 中元素的数乘。于是条件 $\rho(s) = s \otimes e_{\cha}$ 转化为
$$
(\theta(G\times X))(\sigma^* s) = \left(\pr_G^{G \times X}\right)^*\cha \cdot \left(\pr_X^{G \times X}\right)^* s,
$$
这就建立了权空间 $\mathcal{F}(X)_{\cha}$ 的公式。
\end{proof}

\begin{yizhu}[权 $0$ 部分就是不变截面]
权空间分解把「不变截面」精确地放回了它该在的位置：取 $\cha$ 为平凡特征（即 $\cha=1$），条件 $(\theta(G\times X))(\sigma^* s)=(\pr_G)^*\cha\cdot(\pr_X)^*s$ 退化为 $(\theta(G\times X))(\sigma^* s)=(\pr_X)^*s$，这正是 $s$ 沿 $G$-作用不变。因此 $\mathcal{F}(X)_{1}=\mathcal{F}(X)^{G}$，而 §5a 的命题 \ref{prop:InvariantSectionsOfEquivariantInternalHom} 不过是这一分解在内部 Hom 上的特例。反过来说，非平凡权 $\cha\neq1$ 的截面就是「按特征 $\cha$ 变换」的半不变截面；后文 A 侧的单值性正是在这个意义下与权对应。
\end{yizhu}

\begin{prop}[同态空间的权空间分解]
\label{prop:WeightSpaceDecompositionOfHomSpace}
设 $G$ 为对角化群概形，作用在概形 $X$ 上。对任意
$(\mathcal{F}, \theta)$、
$(\mathcal{F}', \theta') \in \Mod_G(X)$，
若
$\mathcal{F}$
有限表现，则从 $\mathcal{F}$ 到 $\mathcal{F}'$ 在 $\Mod(X)$ 中的同态空间有典范的权空间分解
$$
\operatorname{Hom}_{\Mod(X)}(\mathcal{F}, \mathcal{F}') 
= 
\bigoplus_{\cha \in \widehat{G}} \operatorname{Hom}_{\Mod_G(X)}\left( (\mathcal{F}, \theta)(\cha), (\mathcal{F}', \theta') \right).
$$
\end{prop}

\begin{proof}
由于 $\mathcal{F}$ 有限表现，由定义 \ref{defn:InternalHomOfEquivariantModules}，内部 Hom $\mathcal{H} := \mathcal{H}om_{\mathcal{O}_X}(\mathcal{F}, \mathcal{F}')$ 自然带有由 $\alpha \mapsto \theta^{-1} \comp \alpha \comp \theta'$ 给出的 $G$-线性化 $\Theta$，从而 $(\mathcal{H}, \Theta) \in \Mod_G(X)$。

由引理 \ref{lem:GlobalSectionDecomposition}，整体截面空间
$$
\operatorname{Hom}_{\Mod(X)}(\mathcal{F}, \mathcal{F}')
=
\mathcal{H}(X)
$$
% [原文 p.36]
有典范的权空间分解
$
\mathcal{H}(X)
=
\bigoplus_{\cha \in \widehat{G}} \mathcal{H}(X)_{\cha},
$
其中对 $f\in\mathcal{H}(X)$ 与 $\cha\in\widehat{G}$，下列等价成立：
\begin{align*}
f\in \mathcal{H}(X)_{\cha}
\quad
&\Leftrightarrow
\quad
(\Theta(G\times X))(\sigma^* f)
=
\left(\pr_G^{G \times X}\right)^*\cha
\cdot
\left(\pr_X^{G \times X}\right)^* f
\\[2mm]
&\Leftrightarrow
\quad
\theta^{-1} \comp \sigma^* f \comp \theta' = \left(\pr_G^{G \times X}\right)^*\cha \cdot \left(\pr_X^{G \times X}\right)^* f
\\[2mm]
&\Leftrightarrow
\quad
\sigma^* f \comp \theta' = \theta \comp \left( \left(\pr_G^{G \times X}\right)^*\cha \cdot \left(\pr_X^{G \times X}\right)^* f \right)
\\[2mm]
&\Leftrightarrow
\quad
\sigma^* f \comp \theta' = \left( \left(\pr_G^{G \times X}\right)^*\cha \cdot \theta \right) \comp \left(\pr_X^{G \times X}\right)^* f\\[2mm]
&(\because \left(\pr_G^{G \times X}\right)^*\cha \cdot\text{ 是数乘。})
\\[2mm]
&\Leftrightarrow
\quad
\sigma^* f \comp \theta' = \theta(\cha) \comp \left(\pr_X^{G \times X}\right)^* f
\quad(\because \text{$\theta(\cha)$ 的定义})
\end{align*}

该等式恰是 $f$ 成为从 $(\mathcal{F}, \theta(\cha))$ 到 $(\mathcal{F}', \theta')$ 的 $G$-等变态射的条件（定义 \ref{defn:EquivariantSheaves}）。因此 $\mathcal{H}(X)_{\cha} = \operatorname{Hom}_{\Mod_G(X)}\left( (\mathcal{F}, \theta)(\cha), (\mathcal{F}', \theta') \right)$，典范权空间分解得证。
\end{proof}

最后，作为上述计算的推论，我们得到通向主定理证明的一个重要的中间步骤。
\begin{cor}[轨道范畴的嵌入]
\label{thm:EmbeddingOfTheOrbitCategory}
设 $G$ 为对角化群概形，作用在概形 $X$ 上。则遗忘函子 $\Mod_G(X) \rightarrow \Mod(X)$（对象 $(\mathcal{F}, \theta) \mapsto \mathcal{F}$）诱导轨道范畴（orbit category）的全忠实嵌入\footnote{
当 $X$ 局部 Noether 时，对 $\Coh_G(X)$ 与 $\Coh(X)$ 同样成立。
}：
$$
\QCoh_G(X)/\widehat{G}
\hookrightarrow
\QCoh(X).
$$
\end{cor}

\begin{proof}
按定义，轨道范畴中的同态空间为
$$
\operatorname{Hom}_{\QCoh_G(X)/\widehat{G}}\left( (\mathcal{F}, \theta), (\mathcal{F}', \theta') \right)
=
\bigoplus_{\cha \in \widehat{G}} \operatorname{Hom}_{\QCoh_G(X)}\left( (\mathcal{F}, \theta)(\cha), (\mathcal{F}', \theta') \right).
$$
由命题 \ref{prop:WeightSpaceDecompositionOfHomSpace}，该直和恰好与 $\QCoh(X)\subseteq\Mod(X)$ 中的同态空间 $\operatorname{Hom}_{\QCoh(X)}(\mathcal{F}, \mathcal{F}')$ 重合。
% [原文 p.37]
因此，由对象上的 $(\mathcal{F}, \theta) \mapsto \mathcal{F}$ 与态射上的恒等映射定义的函子 $\QCoh_G(X)/\widehat{G} \hookrightarrow \QCoh(X)$ 全忠实。
\end{proof}

